\documentclass[10pt,a4paper]{amsart}
\usepackage[margin=19mm]{geometry}
\usepackage{amsmath,amssymb,mathtools}
\usepackage[scr=rsfs,cal=euler]{mathalfa}
\usepackage{xfrac}
\usepackage[unicode,hidelinks,bookmarks=false]{hyperref}
\newcommand{\ie}{i.e.\ }
\newcommand{\cf}{cf.\ }
\newcommand{\resp}{respectively\ }
\newcommand{\Subsection}{\subsection}
\providecommand{\nobreakdash}{\nobreak\hbox{-}}
\setlength{\emergencystretch}{2em}
\multlinegap0pt
\usepackage{amssymb}
\usepackage{xcolor}
\newtheorem{lemma}{Lemma}
\newtheorem{theorem}{Theorem}
\def\D{\mathcal{D}}
\def\L{\Lambda}
\title{Functional completeness and primitive positive decomposition of relations on finite domains}
\author{Sergiy Koshkin}
\date{Source: arXiv:2606.19492v1 (CC BY 4.0)}
\begin{document}
\maketitle

\noindent\emph{Source and licence.} Functional completeness and primitive positive decomposition of relations on finite domains. Version arXiv:2606.19492v1 (CC BY 4.0), distributed by arXiv under the Creative Commons Attribution 4.0 International licence, http://creativecommons.org/licenses/by/4.0/, as recorded in the arXiv licence metadata for this exact version and shown on the abstract page https://arxiv.org/abs/2606.19492v1. Full licence notice: \texttt{licenses/arxiv-2606.19492-CC-BY-4.0.txt}.\par\smallskip

\noindent\textit{Editorial scope.} Counted unit: the article's theorem FunRed (source0.tex lines 222--227). The article's complete original proof of it is delivered with it (source0.tex lines 228--240, one \texttt{proof} environment of 13 lines). No mathematics is written, repaired or re-derived: the statement and the proof are the article's own lines, in the article's own notation, and no retained line is substituted or re-flowed.  Retained uncounted as the source-local context the proof needs: lines 199--204; lines 211--216.  Nothing was omitted from any retained span, so the package carries no omission record.  No retained line contains a citation, so the wrapper carries no bibliography and no bibliography entry.  No retained line contains a figure environment, an image-inclusion command or drawing code, so no image is delivered and the package declares no asset dependency.  Editorial formatting only, in addition: an AMS \texttt{amsart} preamble that restates the article's own declarations and macros used by the retained lines, namely the environments \texttt{lemma}, \texttt{theorem} on separate counters, together with the standard packages those lines need.  The retained environments print in this document as Lemma 1, Theorem 1 in order of appearance, whereas the article prints the same items with the numbering of its own full document.  The complete native source of the article as distributed in the arXiv e-print for this exact version is bundled unchanged as \texttt{sources/203146/source0.tex}.
\begin{lemma}\label{DomSplit} Any relative $R$ decomposes as
\begin{equation}\label{Fdom}
R(x_1,\dots,x_n)=F(x_1,\dots,x_n)\land\textup{Dom}_R(x_1,\dots,x_{n-1}),
\end{equation}
where $F$ is a multifunction. If $R$ is a partial function then $F$ can be chosen to be a function, i.e. $F(x_1,\dots,x_n)$ if and only if $x_n=f(x_1,\dots,x_{n-1})$.
\end{lemma}

\begin{lemma}\label{MultSel} Let $\D$ be a well-orderable set. Then any $n$-ary multifunction $F$ on $\D$ decomposes as
\begin{equation}\label{MSel}
F(x_1,\dots,x_n)=\exists\,t\big[x_n=f(t,x_1,\dots,x_{n-1})\big],
\end{equation}
where $f:\D^{n}\to\D$ is a function.
\end{lemma}

\begin{theorem}[{\bf Relative reduction}]\label{FunRed} Any $n$-ary relation $R$ on any domain $\D$ decomposes as
\begin{equation}\label{Fred}
R(x_1,\dots,x_n)=\exists\,t\big[x_n=f(t,x_1,\dots,x_{n-1})\big]\land\textup{Dom}_R(x_1,\dots,x_{n-1}),
\end{equation}
where $f:\D^{n}\to\D$ is a function. Moreover, given a functionally complete set, every relation can be decomposed into graphs of functions from this set. In particular, every $n$-ary relation on $\D$ with $n\geq4$ is reducible to ternary relations.  
\end{theorem}

\begin{proof} Formula \eqref{Fred} is a direct consequence of Lemmas \ref{DomSplit} and \ref{MultSel}. To decompose the first conjunct in \eqref{Fred}, set $x_0:=t$ and suppose
$$
f(t,x_1,\dots,x_{n-1})=g\big(h_1(x_{\L_1}),\dots,
h_m(x_{\L_m})\big),
$$
where $\L_i\subseteq\{0,1,\dots,n-1\}$ and $x_{\L}$ is the list of $x_i$ with $i\in\L$. Then
\begin{multline}\label{Fun2Rel}
\exists\,t\big[x_n=f(t,x_1,\dots,x_{n-1})\big]\\ 
= \exists\,t\exists\,t_1\dots\exists\,t_m
\big[x_n=g(t_1,\dots,t_m)\land t_1=h_1(x_{\L_1})\land\dots\land t_m=h_m(x_{\L_m})\big].
\end{multline}
By iterating the procedure, if necessary, we can decompose the first conjunct into graphs of functions from a given complete set. The same process is then applied to $\textup{Dom}_R$. Since the arity is reduced by $1$ at each step the process terminates in finitely many steps. Since there exist functionally complete sets of $2$-input functions, and their graphs are ternary, the last claim follows.
\end{proof}

\end{document}
